\section{Continuation geometry and an exact memory obstruction}
\label{sec:v18-continuation}
\subsection{Normalized continuation laws}
We first separate the transducer from the source to which it is attached.
There are $L$ epochs with finite input alphabets $X_t$ and output alphabets
$Z_t$. A causal kernel $p(z_{1:L}\mid x_{1:L})$ is nonnegative, normalized
for every input word, and has output-prefix marginals independent of later
inputs. The current input can include both a source report and the current
control. It never includes an unavailable parameter, a terminal mark, or
an earlier output not retained by the machine.

Let $\mathcal C_t$, $0\le t\le L$, be the finite-dimensional polytope of
normalized causal kernels from $x_{t+1:L}$ to $z_{t+1:L}$; put
$\mathcal C_L=\{1\}$. For $v\in\mathcal C_t$ define the shift
\[
 (S_{x,z}v)(z_{t+2:L}\mid x_{t+2:L})
   =v(z,z_{t+2:L}\mid x,x_{t+2:L}).
\]
Its total mass $v(z\mid x)$ is independent of the future input word. A
zero shift is allowed.

\begin{definition}[Compatible continuation polytopes]
\label{def:v18-cover}
A continuation cover of width $K=(K_0,\ldots,K_L)$ consists of lists
$V_t=(v_{t,m})_{m\le K_t}$ in $\mathcal C_t$ such that, for every $m,x,z$,
\begin{equation}\label{eq:v18-shift}
 S_{x,z}v_{t,m}\in\operatorname{cone} V_{t+1}.
\end{equation}
It covers $p$ when $p\in\conv V_0$. Lists may repeat generators and may
contain unreachable generators. Thus unused states can be adjoined.
\end{definition}

\begin{theorem}[Continuation realization with a width profile]
\label{thm:v18-continuation}
A causal kernel $p$ has a stochastic transducer realization with an initial
$K_0$-valued register and $K_t$ retained states after epoch $t$ if and only
if it has a compatible continuation cover of width $K$. Initial
randomization is part of the charged initial register. With a deterministic
initial state the condition is $p=v_{0,1}$ and $K_0=1$.

Consequently the set of feasible width profiles depends only on the causal
kernel and its declared input interface, not on a row presentation or a
labeling of the hidden states. Every matrix of conditional continuation
probabilities at cut $t$ has nonnegative rank at most $K_t$. Individual
rank bounds do not replace the simultaneous shift conditions
\eqref{eq:v18-shift}.
\end{theorem}
\begin{proof}
Given a transducer, let $v_{t,m}$ be the suffix response when its register
at time $t$ equals $m$. The remaining rows define this response even for
an unreachable state. If the next row is
$q_{t+1}(z,m'\mid m,x)$, conditioning on its output and new state gives
\[
 S_{x,z}v_{t,m}=\sum_{m'}q_{t+1}(z,m'\mid m,x)v_{t+1,m'}.
\]
The initial state distribution expresses $p$ as a convex combination of
$V_0$.

Conversely choose nonnegative coefficients in every cone membership in
\eqref{eq:v18-shift}. There are only finitely many choices to make. Since
every generator on the right has mass one for each future input word,
the sum of the coefficients for a fixed $m,x,z$ is $v_{t,m}(z\mid x)$.
Summing over $z$ gives one. The coefficients are therefore stochastic
rows $q_{t+1}(z,m'\mid m,x)$. Choose an initial distribution representing
$p$ in $\conv V_0$. Backward induction from $\mathcal C_L=\{1\}$ proves
that the constructed machine has exactly the response $v_{t,m}$ from
each register state. In particular its initial response is $p$. Equality
for open-loop input words implies equality under adaptive input selection:
at each node of an input--output history tree, the same causal conditional
probability is used on the two sides.

For the rank assertion, condition on any positive-probability observed
history at the cut. Its continuation response is a convex combination of
$v_{t,m}$ with weights equal to the posterior register distribution.
Stack these responses as rows, and use all future input--output pairs as
columns. This is a nonnegative factorization through $K_t$ generators.
The converse construction required the same generator families at adjacent
cuts and all their shifts; an unrelated factorization of each cut matrix
provides no such compatibility. The characterization, being stated entirely
in terms of these response polytopes, also proves the asserted invariance.
\end{proof}

Attach such a causal kernel to a marked source experiment. The output law
is obtained by contraction with the source's causal instruments, while its
actual mark is retained. Hence zero private deficiency is equivalently the
existence of a compatible continuation cover whose contraction equals the
target marked behavior in every parameter and feedback row. This is the
geometric form of the factorization in Theorem~\ref{thm:v13-main}; no
conditional dependence on the mark is introduced.

A hidden selector chooses one such cover independently of the source.
It is a convex combination of complete realizations, not a convexification
inside every retained-state simplex at the original width. Its private
implementation uses the disjoint union of the covers and keeps the selector
in the state. A visible selector keeps those fibers in the compared law.
The nonnegative fiber argument in Theorem~\ref{thm:v13-main} therefore
continues to give equality of the visible and private zero sets.

\subsection{An erasure cut}
A simple family makes the width obstruction exact, including its error.
Let $X$ take values in $[N]$ with law $\pi$. Before a specified cut the
source reveals $X$. After the cut it has no reports correlated with $X$.
The controller has no feedback channel carrying $X$ across the cut, and
all earlier emitted reports are prescribed constants. The only permitted
cross-cut information of the simulator is its $K$-valued register and any
independent selector explicitly included in its signature. The terminal
mark is the actual $W=X$. The target's terminal report is $X$.
All memory at other checkpoints is unrestricted.

\begin{theorem}[Exact cut spectrum]
\label{thm:v18-cut}
For this experiment, for private, hidden and visible signatures alike,
\begin{equation}\label{eq:v18-cutvalue}
 \delta_K=1-\sum_{i=1}^{\min(K,N)}\pi_{(i)}.
\end{equation}
For uniform $X$, $\delta_K=(1-K/N)_+$. Thus an erasure cut for $m$
unbiased bits requires $K\ge(1-\epsilon)2^m$ whenever its error is at most
$\epsilon$. The bound holds for every surrounding width profile; equality
is attained when the other checkpoints are unrestricted. If every
$\pi_x>0$, exact transmission has nonnegative cut rank $N$.
\end{theorem}
\begin{proof}
With all post-cut randomness averaged into the decoder, any private
machine gives kernels $e(m\mid x)$ and $d(z\mid m)$. Its joint terminal
law is $\pi_x\sum_m e(m\mid x)d(z\mid m)$. Since the target law is
$\pi_x\1_{z=x}$, total variation is exactly one minus the correct-decoding
probability. Put $u_x=\max_m d(x\mid m)$. Then
\[
 0\le u_x\le1,\qquad \sum_xu_x\le K,\qquad
 \Pp(Z=X)\le\sum_x\pi_xu_x
                  \le\sum_{i=1}^{\min(K,N)}\pi_{(i)}.
\]
The last inequality follows by moving mass in $u$ from a smaller weight
to a larger one until the $K$ largest weights are filled. To attain it,
retain distinct labels for those letters and decode them deterministically;
all other letters are mapped to any retained label. Before the cut one may
use the unrestricted memory to compute this encoding.

For a hidden selector, condition on its value. It is independent of $X$,
so every conditional correct-decoding probability obeys the same bound;
averaging cannot improve it. For a visible selector, variation is the
average of the fiber variations, with the same conclusion. The private
construction works for both signatures. Finally, the exact joint cut
matrix is diagonal with its $N$ positive diagonal entries $\pi_x$.
Its ordinary rank is $N$, so its nonnegative rank is at least $N$ and
its diagonal factorization attains $N$.
\end{proof}

The absence of an auxiliary cross-cut feedback channel is essential. An
unpriced controller storing $X$ and returning it as a later control would
change the experiment, not disprove the cut bound. Likewise a selector
correlated with $X$ is not an independent hidden selector. These distinctions
are the same ones needed in the collision certificate below.
